% LSPV-3 REPORT::LOGCOF1::2026-10-12T19:32:31Z
\documentclass[11pt]{article}
\usepackage[margin=1in]{geometry}
\usepackage{amsmath,amssymb,booktabs,array,hyperref}
\usepackage[T1]{fontenc}
\hypersetup{colorlinks=true,urlcolor=blue,linkcolor=blue,citecolor=blue}
\newcommand{\Status}[1]{\textbf{#1}}
\newcommand{\SAT}{\mathrm{SAT}}
\title{LSPV-3 LOGCOF1: Logarithmic-width cofactor tables defeat late selection as a universal bottleneck}
\author{UNIVERSAL\_CORE\_REDTEAM (LSPV-3)}
\date{2026-10-12T19:32:31Z}

\begin{document}
\maketitle

\section*{Executive status}
\Status{PROVED/EXACT:} if a deterministic random-access decider \(A\) exists,
then for any fixed \(c\) there is an equivalent wrapper that withholds
\(k(N)\le c\log_2 N\) source bits, evaluates \(A\) on all \(2^k\) completions,
and performs only \(O(1)\) work after the final source read. Its fully charged
cost is
\[
 T_B(N)\le 2^kT_A(N)+O(N+2^k),\qquad
 S_B(N)\le S_A(N)+N+2^k+O(\log N).
\]
Thus polynomially many unresolved cofactors do not force a superpolynomial
late decision-selection phase. \Status{FALSIFIED/RETIRED:}
\texttt{POLYNOMIAL-COFACTOR-AMBIGUITY-AS-SELECTION-LB1}.

The theorem relocates and duplicates work; it does not make total SAT decision
cheap, improve an exact-SAT exponent, or prove any total-time lower bound.
\[
 \boxed{\text{P versus NP remains open.}}
\]

\section{Predecessor write-recovery preflight}
The predecessor LSPV-2 report
\path{research/agents/lspv-2/2026-10-13_0423_LSPV-2_COVERCOF1.tex}
exists canonically, is substantive, and has the required header and sentinel.
Its experiment and result artifacts also exist. However root \path{main.tex}
still contains the older LSPV-2 region and therefore the predecessor run has
not completed its mandatory same-run root reconciliation. This is classified
as \Status{PROMPT / PROCEDURE INTERRUPTION}, not missing science. A
body-preserving recovery request was issued; no duplicate predecessor report
was created and this red-team run continued.

\section{CORPUS REVIEWED}
The run enumerated the current repository directories and read root
\path{main.tex}, \path{AGENTS.md}, \path{CODEX.md},
\path{.codex/skills/pnp-research/SKILL.md}, the current task specifications,
actual \path{main4.tex} and \path{main5.tex}, and the newest verified MAIN
snapshot REPPROMISE1. It read the materially relevant newest specialists:
HIERFIRE1, COVERCOF1, REVUNCOMP1/CAUSALCANCEL1, and READFLOOR1; and searched
historical LSPV reports and explorations for cofactor, decision-selection,
late-read, query-floor, wrapper-orbit, trace, source-copy, product-streaming,
HEC, VSC-HEC, G-PPC-HEC, CRS-PPC, CCSPC and RCCN routes. Governance and cited
experiment artifacts were included. \path{tasks/TASK-008_SLICE-UNSAT.md}
remains absent and was not reconstructed.

\section{Primary-literature and transfer preflight}
The two obstacles are (i) making decision selection an implementation-invariant
event, and (ii) transferring a lower bound from query, circuit, branching
program, VLSI, communication, or proof models to arbitrary uniform
random-access computation. Shannon/Boole cofactor expansion is an exact
Boolean identity and supports the upper-bound construction, not hardness
\cite{Shannon1949}. Rivest--Vuillemin query lower bounds remain exact only in
their decision-tree model \cite{RivestVuillemin1976}. Pippenger--Fischer
oblivious simulation is an upper-bound transformation with overhead, not a
theorem forcing a chosen internal phase \cite{PippengerFischer1979}.
Relativization and algebrization continue to block a generic black-box
promotion \cite{BGS1975,AW2009}. These are direct model statements; importing
their vocabulary without a compiler is a failed transfer.

\section{Exact target attacked}
For a run of a decider \(M\) on an \(N\)-bit string, let
\(\tau_M(x)\) be the first time all real source positions have been read and
let \(R_M(x)=T_M(x)-\tau_M(x)\). COVERCOF1 already proves that a single
withheld bit permits \(R_M(x)=O(1)\) after two virtual runs. LOGCOF1 attacks the
proposed repair that many unresolved final choices restore a lower bound:
\[
 \texttt{POLYNOMIAL-COFACTOR-AMBIGUITY-AS-SELECTION-LB1}:
 \quad 2^{k(N)}=N^{O(1)}
 \Longrightarrow \text{unavoidable superpolynomial final selection work}.
\]
The implication is false for every decidable language, including any explicit
3-SAT clean core.

\section{LOG-WIDTH-COFACTOR-NORMAL-FORM1}
\paragraph{Theorem (\Status{PROVED/EXACT}).}
Let \(A\) be any total uniform deterministic random-access decider for a
language \(L\). Let \(W_N\subseteq[N]\) be a uniformly computable set of
\(k(N)\) withheld coordinates. There is a total uniform deterministic decider
\(B\) for \(L\) that:
\begin{enumerate}
 \item reads and stores all coordinates outside \(W_N\);
 \item simulates \(A\) on every completion \(u\in\{0,1\}^{k(N)}\), storing the
 one-bit outputs in a \(2^{k(N)}\)-entry table;
 \item reads the real withheld coordinates while accumulating their table
 index; and
 \item after reading the last real coordinate, performs one random-access
 table lookup and outputs that bit.
\end{enumerate}
In a word-RAM with the same random-access convention as \(A\),
\[
 R_B(x)=O(1),\quad
 T_B(N)\le 2^kT_A(N)+O(N+2^k),\quad
 S_B(N)\le S_A(N)+N+2^k+O(\log N).
\]
If \(k(N)\le c\log_2N\) for fixed \(c\), polynomial time and space are
preserved.

\paragraph{Proof.}
For each real input \(x\), exactly one completion equals \(x\), namely the
entry indexed by \(x|_{W_N}\). Totality of \(A\) makes every virtual run halt.
The wrapper answers virtual source queries from its stored nonwithheld bits and
the current hardwired completion. The real withheld bits are read only after
all table entries exist. The index is accumulated online, so the last source
read is followed by one table lookup and output. Summing the \(2^k\)
simulations and charging the source store and table gives the displayed bounds.
When \(2^k\le N^c\), multiplication by the arbitrary fixed polynomial of
\(A\) remains polynomial; its exponent is increased, never erased.
\(\square\)

\paragraph{P=NP counterfactual.}
Assume \(\SAT\in\mathrm P\) via a decider \(A\) of unknown exponent \(K_A\).
For fixed \(c\), LOGCOF1 gives a uniform SAT decider in time
\(N^{K_A+c+O(1)}\) whose final selection among \(N^c\) cofactors is constant
time. This is fully consistent with \(\mathrm P=\mathrm{NP}\). Therefore a
contradiction cannot come from polynomial cofactor ambiguity or the location
of final selection. Banning the wrapper would assume the desired universal
obligation rather than derive it.

\section{Semantic compatibility and hostile controls}
For CNF \(F\) and partial assignment \(\alpha\), define
\[
 H(F,\alpha)=\mathbf 1[\exists z\supseteq\alpha\;z\models F].
\]
The table contains only terminal language bits of complete source strings. It
does not expose a witness, a compatibility relation over assignments, a
separator, proof, communication transcript, or reconstruction. The table may
contain polynomially many bits even when the semantic solution set is
exponential; conversely it may have high redundancy on easy families.

The identity is unaffected by Horn forward chaining, affine/Tseitin Gaussian
elimination, 2-SAT implication graphs, bounded-treewidth dynamic programming,
parity, independent products, padding, redundant clauses, bit reversal,
complement encodings, compact circuit/branching-program presentations, or
DIRECT/COPY/DEAD/DELAY/INTERP/RECOMP wrappers. These controls show that
polynomial ambiguity is not a hardness certificate; they do not make general
3-SAT easy.

\section{Bounded reproducible audit}
\path{experiments/review/log_width_cofactor_selection_audit.py} tested sizes
\(1\le N\le11\), all inputs, nine predicate families, three encodings,
prefix/suffix/spaced withheld sets, every \(0\le k\le\lfloor\log_2N\rfloor\),
and six wrapper labels. It performed 6,811,182 assertions in 12.2071 seconds,
with zero failures. The largest bounded table had eight entries.
\[
\begin{array}{ll}
\text{script SHA-256}&
\texttt{e9ddcd9a884673413f8f4bbadcd21f8f56ef352c5ba3334209c5c2e37b0b7470}\\
\text{result SHA-256}&
\texttt{66268b16b717cbcdd66c8b1004f007958a0771eca2073e1426ad448f54d5e79e}
\end{array}
\]
The invocation completed below 30 seconds, so no chunking was required.
\Status{COMPUTATIONAL:} the audit is a finite identity regression and no
asymptotic runtime is inferred from it.

\section{DUPLICATION / CONTRADICTION AUDIT}
COVERCOF1 is the immediate one-bit-width theorem; LOGCOF1 is its new
logarithmic-width closure and quantifies the exact polynomial exponent charge.
PRODSTREAM1 separately collapses product answer vectors to a streaming
accumulator. READFLOOR1 proves a near-linear source-probe floor but explicitly
does not force post-read computation. READCEIL1 gives the matching linear
query ceiling. FIELDMASK1, CAUSALCANCEL1 and REVUNCOMP1 retire fixed trace
masks, causal liveness and logical erasure. LOGCOF1 does not duplicate those
results and creates no contradiction: it sharpens their common implementation-
location guard.

\section{WHAT CHANGED RELATIVE TO ROOT MAIN}
Root \path{main.tex} still lists \texttt{DECISION-SELECTION-WORK1} as open
without the logarithmic cofactor-table guard. The LSPV-3 region should now
record that any viable decision-selection target must charge the full
precomputation and cannot rely on polynomially many final alternatives,
late source availability, or a superpolynomial suffix after complete coverage.
The direct total-time clean-core lower bound remains open.

\section{ADOPTED / REJECTED / RETIRED PRIOR COMPONENTS}
\begin{center}
\begin{tabular}{>{\raggedright}p{0.22\linewidth}p{0.68\linewidth}}
\toprule
\Status{ADOPTED} & SAME\_REALIZATION, arbitrary exponent accounting, one-bit
output, totality, wrapper tests, and READFLOOR1 source sensitivity.\\
\Status{ADOPTED} & COVERCOF1 as the \(k=1\) base case.\\
\Status{REJECTED} & Inferring runtime from semantic-state count, number of
cofactors, table width, or the chronological location of source reads.\\
\Status{FALSIFIED/RETIRED} &
\texttt{POLYNOMIAL-COFACTOR-AMBIGUITY-AS-SELECTION-LB1}.\\
\bottomrule
\end{tabular}
\end{center}

\section{STALE OR OVERSTATED CLAIMS FOUND}
Any claim that the solver must do superpolynomial work after learning which of
polynomially many source completions is real is overstated. A wrapper can do
all charged decision work before learning that index. The safe surviving claim
is only that total work might be superpolynomial; proving that is precisely the
open problem.

\section{Barrier-transfer and post-quantum ledger}
\begin{itemize}
 \item Cofactor expansion: \Status{EXACT OBJECT/THEOREM TRANSFER}; it gives an
 upper-bound normal form, not hardness.
 \item Decision-tree and branching-program lower bounds:
 \Status{DIRECTLY TRANSFERABLE} only inside their declared models; promotion to
 arbitrary computation without a compiler is \Status{FAILED TRANSFER}.
 \item Communication cuts, VLSI routing, pebbling, observability and
 controllability: \Status{ANALOGY/CANDIDATE}; the missing all-decider compiler
 must preserve every cost.
 \item Schur squares, quadratic relations, Hilbert functions, syzygies and
 graded Betti numbers: \Status{FAILED TRANSFER} for this target. Hidden
 algebraic structure may accelerate the algorithm.
 \item Collision entropy is exactly computable for the table map, but high
 collision or small image supplies no decision lower bound:
 \Status{FAILED TRANSFER}.
 \item Worst-case/average-case and special-soundness extraction:
 \Status{ANALOGY/CANDIDATE}; no unconditional black-box extractor follows.
\end{itemize}
The normal form relativizes and is a Boolean cofactor identity, so it also
survives algebraic extensions. It supplies no large constructive circuit
property and does not evade Natural Proofs. It is a guard, not a barrier escape.

\section{Strict status and smallest repair}
\begin{center}
\begin{tabular}{lp{0.67\linewidth}}
\toprule
\Status{PROVED/EXACT} & LOG-WIDTH-COFACTOR-NORMAL-FORM1 and all displayed
costs.\\
\Status{COMPUTATIONAL} & 6,811,182 bounded assertions; zero failures.\\
\Status{FALSIFIED/RETIRED} & polynomial cofactor ambiguity as unavoidable
late selection work.\\
\Status{OPEN} & a wrapper-invariant total-cost observable or direct
superpolynomial lower bound for every decider of the fixed NP-complete core.\\
\Status{OPEN} & a model-matched unconditional contradiction.\\
\bottomrule
\end{tabular}
\end{center}
The smallest repair to \texttt{DECISION-SELECTION-WORK1} is: charge the whole
execution, including all virtual runs and table storage, and require a
source-independent lower bound invariant under cofactor precomputation. Any
late-phase formulation must exclude the normal form by a restriction that is
itself transferred from every ordinary SAT decider. That missing theorem is
separation-strength and remains open.

\begin{thebibliography}{9}
\bibitem{Shannon1949}
C.~E. Shannon, The Synthesis of Two-Terminal Switching Circuits,
\emph{Bell System Technical Journal} 28(1), 1949.
\href{https://doi.org/10.1002/j.1538-7305.1949.tb03624.x}{doi:10.1002/j.1538-7305.1949.tb03624.x}.
\bibitem{RivestVuillemin1976}
R.~L. Rivest and J.~Vuillemin, On Recognizing Graph Properties from
Adjacency Matrices, \emph{Theoretical Computer Science} 3(3), 1976.
\href{https://doi.org/10.1016/0304-3975(76)90053-0}{doi:10.1016/0304-3975(76)90053-0}.
\bibitem{PippengerFischer1979}
N.~Pippenger and M.~J. Fischer, Relations Among Complexity Measures,
\emph{JACM} 26(2), 1979.
\href{https://doi.org/10.1145/322123.322138}{doi:10.1145/322123.322138}.
\bibitem{BGS1975}
T.~Baker, J.~Gill, and R.~Solovay, Relativizations of the P=?NP Question,
\emph{SIAM Journal on Computing} 4(4), 1975.
\href{https://doi.org/10.1137/0204037}{doi:10.1137/0204037}.
\bibitem{AW2009}
S.~Aaronson and A.~Wigderson, Algebrization: A New Barrier in Complexity
Theory, \emph{TOCT} 1(1), 2009.
\href{https://doi.org/10.1145/1490270.1490272}{doi:10.1145/1490270.1490272}.
\end{thebibliography}

% END_OF_VERIFIED_REPORT::LSPV-3::LOGCOF1
\end{document}
